% ============================================================
% PHỤ LỤC
% ============================================================
\appendix
\chapter{Phụ lục}
\label{chap:appendix}

\section{Cấu trúc mã nguồn dự án}
\label{sec:code_structure}

\begin{lstlisting}[language={}, caption={Cấu trúc thư mục mã nguồn}, label={lst:project_structure}]
speech_emo_recognition/
|-- IEMOCAP_logspec200.pkl        # Dataset dac trung Log Spectrogram
|-- train_ser.py                  # Script huan luyen chinh
|-- model.py                     # Dinh nghia tat ca kien truc mang
|-- data_utils.py                # Xu ly va tai du lieu
|-- run_experiments.py           # Chay tu dong tat ca thuc nghiem
|-- crossval_efficientnet.py     # 5-Fold Cross-Validation
|-- ablation_study.py            # Ablation study
|-- plot_comparison.py           # Ve bieu do so sanh tong hop
|-- test_demo.py                 # Demo nhan dang cam xuc tu file .wav
|-- features_extraction/         # Scripts trich xuat dac trung
|   |-- extract_features.py      # Script trich xuat chinh
|   |-- features_util.py         # Ham tien ich xu ly tin hieu
|   |-- database.py              # Giao dien bo du lieu IEMOCAP
|   `-- noise_wav/               # File nhieu de augmentation
|-- requirements.txt             # Danh sach thu vien Python
`-- *.pth                        # Trong so cac mo hinh da huan luyen
\end{lstlisting}

\section{Hướng dẫn sử dụng}
\label{sec:usage_guide}

\subsection{Cài đặt môi trường}

\begin{lstlisting}[language=bash, caption={Cài đặt môi trường Python}]
# Tao moi truong Conda
conda create -n pt310 python=3.10
conda activate pt310

# Cai dat dependencies
pip install -r requirements.txt
\end{lstlisting}

\subsection{Huấn luyện mô hình}

\begin{lstlisting}[language=bash, caption={Lệnh huấn luyện EfficientNet-B0 (mô hình tốt nhất)}]
python train_ser.py IEMOCAP_logspec200.pkl \
    --ser_model efficientnet \
    --val_id 1F --test_id 1M \
    --num_epochs 60 --batch_size 64 --lr 0.0001 \
    --seed 100 --gpu 1 \
    --save_label efficientnet_focal_loss \
    --loss_type focal \
    --emix emix-s \
    --pretrained
\end{lstlisting}

\subsection{5-Fold Cross-Validation}

\begin{lstlisting}[language=bash, caption={Chạy 5-Fold Cross-Validation}]
python crossval_efficientnet.py
\end{lstlisting}

\subsection{Ablation Study}

\begin{lstlisting}[language=bash, caption={Chạy Ablation Study}]
python ablation_study.py
\end{lstlisting}

\subsection{Demo nhận dạng cảm xúc}

\begin{lstlisting}[language=bash, caption={Demo nhận dạng cảm xúc từ file .wav}]
python test_demo.py \
    --model efficientnet \
    --pth efficientnet_focal_loss.pth \
    --wav <duong_dan_file_wav>
\end{lstlisting}

\subsection{Vẽ biểu đồ so sánh}

\begin{lstlisting}[language=bash, caption={Tạo biểu đồ so sánh tổng hợp}]
python -X utf8 plot_comparison.py
\end{lstlisting}

Các biểu đồ đầu ra bao gồm tệp \texttt{comparison\_wa\_ua.png} thể hiện biểu đồ cột nhóm so sánh WA/UA; tệp \texttt{comparison\_per\_class.png} hiển thị bản đồ nhiệt (heatmap) độ chính xác theo từng lớp; tệp \texttt{comparison\_params\_vs\_wa.png} biểu diễn biểu đồ phân tán mối quan hệ giữa số lượng tham số và WA; tệp \texttt{summary\_figure.png} biểu diễn biểu đồ tổng hợp chất lượng cao; và cuối cùng là tệp \texttt{crossval\_results.png} chứa kết quả trực quan hóa của quy trình đánh giá chéo 5-Fold.


\section{Tham số dòng lệnh huấn luyện}
\label{sec:cli_params}

\begin{table}[H]
\centering
\caption{Danh sách tham số dòng lệnh của \texttt{train\_ser.py}}
\label{tab:cli_params}
\begin{tabularx}{\textwidth}{p{3.5cm}p{3.2cm}X}
\toprule
\textbf{Tham số} & \textbf{Mặc định} & \textbf{Mô tả} \\
\midrule
\texttt{features\_file} & (bắt buộc) & Đường dẫn file đặc trưng \texttt{.pkl} \\
\texttt{--ser\_model} & \texttt{ser\_alexnet} & Kiến trúc mô hình: \texttt{alexnet}, \texttt{alexnet\_gap}, \texttt{resnet}, \texttt{densenet}, \texttt{efficientnet}, \texttt{mobilenet} \\
\texttt{--val\_id} & \texttt{1F} & ID speaker làm validation \\
\texttt{--test\_id} & \texttt{1M} & ID speaker làm test \\
\texttt{--num\_epochs} & 200 & Số epochs huấn luyện \\
\texttt{--batch\_size} & 32 & Kích thước mini-batch \\
\texttt{--lr} & 0.0001 & Tốc độ học \\
\texttt{--seed} & 100 & Random seed \\
\texttt{--gpu} & 1 & Sử dụng GPU (1=có, 0=không) \\
\texttt{--save\_label} & None & Nhãn lưu mô hình \\
\texttt{--loss\_type} & \texttt{ce} & Hàm lỗi: \texttt{ce} (Cross Entropy) hoặc \texttt{focal} (Focal Loss) \\
\texttt{--emix} & None & EMix augmentation: \texttt{emix-s}, \texttt{emix-n}, \texttt{emix-ns} \\
\texttt{--pretrained} & False & Sử dụng ImageNet pre-trained weights \\
\texttt{--oversampling} & False & Random oversampling \\
\texttt{--mixup} & False & MixUp augmentation \\
\texttt{--load\_model} & None & Nạp trọng số mô hình sẵn có \\
\bottomrule
\end{tabularx}
\end{table}


\section{Danh sách file trọng số mô hình}
\label{sec:model_weights}

\begin{table}[H]
\centering
\caption{Danh sách file trọng số mô hình đã huấn luyện}
\label{tab:model_weights}
\begin{tabularx}{\textwidth}{l X r}
\toprule
\textbf{File} & \textbf{Mô tả} & \textbf{Kích thước} \\
\midrule
\texttt{alexnet\_baseline.pth} & AlexNet Baseline (CE) & 228 MB \\
\texttt{alexnet\_final.pth} & AlexNet Fine-tuned & 228 MB \\
\texttt{alexnet\_gap\_cross\_entropy.pth} & AlexNet GAP (CE) & 9.9 MB \\
\texttt{alexnet\_gap\_focal\_loss.pth} & AlexNet GAP (Focal) & 9.9 MB \\
\texttt{resnet\_focal\_loss.pth} & ResNet-18 (Focal + EMix-S) & 44.8 MB \\
\texttt{densenet\_focal\_loss.pth} & DenseNet-121 (Focal + EMix-S) & 28.5 MB \\
\texttt{efficientnet\_focal\_loss.pth} & EfficientNet-B0 (Focal + EMix-S) & 16.4 MB \\
\texttt{mobilenet\_focal\_loss.pth} & MobileNet-V2 (Focal + EMix-S) & 9.2 MB \\
\texttt{efficientnet\_focal\_cv\_fold*.pth} & 5-Fold CV weights (5 files) & 16.4 MB mỗi file \\
\bottomrule
\end{tabularx}
\end{table}
